\documentclass[../../main.tex]{subfiles}% 注意这里的文件路径不能用 ./main.tex ,否则用latexmk编译子文件会报错
\graphicspath{{\subfix{./image/}}} % 指定图片目录,后续可以直接使用图片文件名
% 注意这里的文件路径不能用 ../../image/ ,否则用latexmk编译子文件会报错

% 例如:
% \begin{figure}[H]
% \centering
% \includegraphics[scale=0.3]{图.png}
% \caption{}
% \label{figure:图}
% \end{figure}
% 注意:上述\label{}一定要放在\caption{}之后,否则引用图片序号会只会显示??.

\begin{document}

\section{主理想整环上有限生成模的应用}

\begin{theorem}
设$G$是有限生成的Abel群，群运算为加法，则存在整数$n\geqslant0$及整数组$d_1,d_2,\dots,d_k$满足
\begin{enumerate}[(1)]
\item $d_i\mid d_{i+1},\ 1\leqslant i\leqslant k-1$;
\item $G\cong \underbrace{\mathbb{Z}\oplus\mathbb{Z}\oplus\cdots\oplus\mathbb{Z}}_{n\text{个}}\oplus\mathbb{Z}_{d_1}\oplus\cdots\oplus\mathbb{Z}_{d_k}$.
\end{enumerate}
\end{theorem}
\begin{proof}
因$G$是有限生成的，故$G$是有限生成的$\mathbb{Z}$-模。又$\mathbb{Z}$是主理想整环，故
\begin{align*}
G\cong \mathbb{Z}x_1\oplus\mathbb{Z}x_2\oplus\cdots\oplus\mathbb{Z}x_n\oplus\mathbb{Z}y_1\oplus\mathbb{Z}y_2\oplus\cdots\oplus\mathbb{Z}y_k,
\end{align*}
且
\begin{align*}
\mathrm{ann}\,x_i=\langle0\rangle,\ 1\leqslant i\leqslant n,\quad \mathrm{ann}\,y_j=\langle d_j\rangle,\ d_j\mid d_{j+1},\ 1\leqslant j\leqslant k-1.
\end{align*}
当$\mathrm{ann}\,x_i=\langle0\rangle$时，$\mathbb{Z}x_i\cong\mathbb{Z}$；而$\mathbb{Z}y_j\cong\mathbb{Z}_{d_j}$，故定理成立.

\end{proof}
\begin{remark}
$n$与$d_1,d_2,\dots,d_k$由$G$唯一确定。$n$是$G$作为$\mathbb{Z}$-模的秩，称为$G$的\textbf{Betti数}；$d_1,d_2,\dots,d_k$称为$G$的\textbf{扭系数}。有限生成的Abel群在代数学、拓扑学中均有重要应用。
\end{remark}

\begin{theorem}
设$G$为有限生成的Abel群，则
\begin{align*}
G\cong \bigoplus_{i,j}\mathbb{Z}_{d_{i,j}}\oplus\underbrace{\mathbb{Z}\oplus\mathbb{Z}\oplus\cdots\oplus\mathbb{Z}}_{n\text{个}},
\end{align*}
其中$d_{i,j}=p_i^{k_{i,j}}$，$p_i$为素数。
\end{theorem}
\begin{remark}
$d_{i,j}=p_i^{k_{i,j}}$称为$G$的\textbf{初等因子}。
\end{remark}

\begin{theorem}
设$G,G_1$都是有限生成的Abel群，则下列条件等价：
\begin{enumerate}[(1)]
\item $G$与$G_1$同构；
\item $G$与$G_1$有相同的Betti数与扭系数；
\item $G$与$G_1$有相同的Betti数及初等因子。
\end{enumerate}
\end{theorem}
\begin{remark}
上述三个定理合称为\textbf{有限生成的Abel群的基本定理}。
\end{remark}

\begin{theorem}
设$\mathscr{A}$是域$F$上线性空间$V$的线性变换，由$\mathscr{A}$定义的$F[\lambda]$-模$V$的不变因子为$d(\lambda)$，即$V=F[\lambda]z,\ \mathrm{ann}\,z=\langle d(\lambda)\rangle$，则
\begin{enumerate}[(1)]
\item $\deg d(\lambda)=\dim V$，且$z,\mathscr{A}z,\dots,\mathscr{A}^{n-1}z$为$V$的一组基；
\item 设$d(\lambda)=\lambda^n+a_{n-1}\lambda^{n-1}+\cdots+a_1\lambda+a_0$，则$\mathscr{A}$在上述基下的矩阵为
\begin{align*}
M(\mathscr{A};z,\mathscr{A}z,\dots,\mathscr{A}^{n-1}z)=
\begin{pmatrix}
0&&&-a_0\\
1&0&&-a_1\\
&\ddots&&\vdots\\
&&0&-a_{n-2}\\
&&1&-a_{n-1}
\end{pmatrix}.
\end{align*}
\end{enumerate}
\end{theorem}
\begin{proof}
由$\mathrm{ann}\,z=\langle d(\lambda)\rangle$，$f(\lambda)z=0$当且仅当$d(\lambda)\mid f(\lambda)$，故$\deg f(\lambda)\geqslant\deg d(\lambda)=n$，因此$z,\mathscr{A}z,\dots,\mathscr{A}^{n-1}z$线性无关。
任取$x\in V$，由$V=F[\lambda]z$，存在$f(\lambda)\in F[\lambda]$使得$x=f(\lambda)z$。作带余除法$f(\lambda)=d(\lambda)\cdot q(\lambda)+r(\lambda)$，其中$r(\lambda)=0$或$\deg r(\lambda)<\deg d(\lambda)$，则$x=r(\lambda)z\in L(z,\mathscr{A}z,\dots,\mathscr{A}^{n-1}z)$，故(1)成立。

对$0\leqslant k\leqslant n-2$，有$\mathscr{A}(\mathscr{A}^kz)=\mathscr{A}^{k+1}z$；
\begin{align*}
\mathscr{A}(\mathscr{A}^{n-1}z)=\mathscr{A}^nz=\left(d(\lambda)-\sum_{k=0}^{n-1}a_k\lambda^k\right)z=-\sum_{k=0}^{n-1}a_k\mathscr{A}^kz,
\end{align*}
故(2)成立.

\end{proof}
\begin{remark}
上述矩阵称为$d(\lambda)$的\textbf{相伴矩阵}；满足该定理的线性空间$V$称为$\mathscr{A}$的\textbf{循环空间}。此时$\mathscr{A}$的特征多项式与极小多项式相等。
\end{remark}

\begin{theorem}
设$\mathscr{A}$是域$F$上线性空间$V$的线性变换，由$\mathscr{A}$定义的$F[\lambda]$-模$V$的不变因子为$d_1(\lambda),d_2(\lambda),\dots,d_s(\lambda)$，满足$d_i(\lambda)\mid d_{i+1}(\lambda),\ 1\leqslant i\leqslant s-1$，则
\begin{enumerate}[(1)]
\item $V$中存在一组基，使得$\mathscr{A}$在该基下的矩阵为准对角阵$\mathrm{diag}(B_1,B_2,\dots,B_s)$，其中$B_i$是$d_i(\lambda)$的相伴矩阵；
\item $\displaystyle\sum_{i=1}^s\deg d_i(\lambda)=\dim V$;
\item $\mathrm{ann}\,V=\langle d_s(\lambda)\rangle$，即$\mathscr{A}$的极小多项式是$d_s(\lambda)$；
\item $\mathscr{A}$的特征多项式为$f(\lambda)=\det(\lambda\mathrm{id}-\mathscr{A})=\displaystyle\prod_{i=1}^s d_i(\lambda)$;
\item $f(\mathscr{A})=0$.
\end{enumerate}
\end{theorem}
\begin{proof}
(1) $V$作为$F[\lambda]$-模有循环子模分解
\begin{align*}
V=\bigoplus_{i=1}^s F[\lambda]z_i,
\end{align*}
其中$\mathrm{ann}\,z_i=\langle d_i(\lambda)\rangle,\ d_i(\lambda)\mid d_{i+1}(\lambda),\ i=1,2,\dots,s-1$。
令$V_i=F[\lambda]z_i$，则$V_i$是$\mathscr{A}$的不变子空间。记$\mathscr{A}_i=\mathscr{A}|_{V_i}$，由上一定理，$z_i,\mathscr{A}z_i,\dots,\mathscr{A}^{n_i-1}z_i$（$n_i=\deg d_i(\lambda)$）是$V_i$的基，$\mathscr{A}_i$在该基下的矩阵为$d_i(\lambda)$的相伴矩阵。结合$V=V_1\oplus V_2\oplus\cdots\oplus V_s$，结论成立。

(2) 
\begin{align*}
\dim V=\sum_{i=1}^s\dim V_i=\sum_{i=1}^s\deg d_i(\lambda).
\end{align*}

(3) 设$\mathrm{ann}\,V=\langle d(\lambda)\rangle$，则$d(\lambda)z_i=0$，故$d(\lambda)\in\mathrm{ann}\,z_i=\langle d_s(\lambda)\rangle$，即$d_s(\lambda)\mid d(\lambda)$。
又$d_i(\lambda)\mid d_s(\lambda)\ (1\leqslant i\leqslant s)$，故$d_s(\lambda)z_i=0\ (\forall\ 1\leqslant i\leqslant s)$，即$d_s(\lambda)x=0\ (\forall x\in V)$，因此$d(\lambda)\mid d_s(\lambda)$，故$d(\lambda)=d_s(\lambda)$，即$\mathscr{A}$的极小多项式为$d_s(\lambda)$。

(4)
\begin{align*}
f(\lambda)=\det(\lambda\mathrm{id}-\mathscr{A})=\prod_{i=1}^s\det(\lambda\mathrm{id}_{V_i}-B_i)=\prod_{i=1}^s d_i(\lambda).
\end{align*}

(5)
\begin{align*}
f(\mathscr{A})=\prod_{i=1}^s d_i(\mathscr{A})=\left(\prod_{i=1}^{s-1}d_i(\mathscr{A})\right)d_s(\mathscr{A})=0.
\end{align*}

\end{proof}
\begin{remark}
该定理给出的矩阵称为$\mathscr{A}$的\textbf{有理标准形}。两个$n$阶方阵或线性变换相似当且仅当它们有相同的有理标准形。
\end{remark}

\begin{theorem}
设$F$是代数封闭域，$V$是$F$上有限维线性空间，$\mathscr{A}$是$V$的线性变换，则$V$中存在一组基，使得$\mathscr{A}$在该基下的矩阵为准对角阵$\mathrm{diag}(C_1,C_2,\dots,C_t)$，其中
\begin{align*}
C_i=
\begin{pmatrix}
\lambda_i&&&\\
1&\lambda_i&&\\
&\ddots&\ddots&\\
&&1&\lambda_i
\end{pmatrix},\quad i=1,2,\dots,t.
\end{align*}
\end{theorem}
\begin{proof}
由主理想整环上模的分解定理，$F[\lambda]$-模$V$可分解为循环子模的直和
\begin{align*}
V=\bigoplus_{i=1}^t F[\lambda]y_i,\quad \mathrm{ann}\,y_i=\langle(\lambda-\lambda_i)^{r_i}\rangle,\ i=1,2,\dots,t.
\end{align*}
令$V_i=F[\lambda]y_i$，则$V_i$是$\mathscr{A}$的不变子空间，且$y_i,(\mathscr{A}-\lambda_i\mathrm{id})y_i,\dots,(\mathscr{A}-\lambda_i\mathrm{id})^{r_i-1}y_i$是$V_i$的一组基。

对$1\leqslant k_i\leqslant r_i-2$，有
\begin{align*}
\mathscr{A}(\mathscr{A}-\lambda_i\mathrm{id})^{k_i}y_i=\lambda_i(\mathscr{A}-\lambda_i\mathrm{id})^{k_i}y_i+(\mathscr{A}-\lambda_i\mathrm{id})^{k_i+1}y_i,
\end{align*}
且
\begin{align*}
\mathscr{A}(\mathscr{A}-\lambda_i\mathrm{id})^{r_i-1}y_i=\lambda_i(\mathscr{A}-\lambda_i\mathrm{id})^{r_i-1}y_i.
\end{align*}
故$\mathscr{A}|_{V_i}$在上述基下的矩阵为$C_i$。结合$V=V_1\oplus V_2\oplus\cdots\oplus V_t$，存在$V$的基使得$\mathscr{A}$的矩阵为Jordan标准形.

\end{proof}
\begin{remark}
$(\lambda-\lambda_i)^{r_i}$称为$\mathscr{A}$的\textbf{初等因子}。$\mathscr{A}$的Jordan标准形除Jordan块$C_i$的次序外唯一确定；两个矩阵或线性变换相似当且仅当它们有相同的Jordan标准形。求解有理标准形或Jordan标准形的核心是求不变因子或初等因子。
\end{remark}



\end{document}